\clearpage
\fancyhead[L]{\small HALLASCHKA. ELEMENTA ECLIPSIUM.}
\SourceLeaf{624}
\begin{center}\small Göttingische gelehrte Anzeigen. Stück 17. Seiten 167..168. 29. Januar 1820.\\[3pt]\rule{.72\linewidth}{.3pt}\end{center}

\emph{Elementa eclipsium quas patitur tellus, luna eam inter et solem versante, ab a. {\upshape 1816} usque ad a. {\upshape 1860}, ex tabulis astronomicis recentissime conditis et calculo parallactico deducta, typo ecliptico et tabulis projectionis geographicis collustraia, a} \textsc{Cassiano Hallaschka} Prag 1816. Typis Theophili Haase. 107 Seiten in Quart, dazu 20 Kupfertafeln.

\begingroup\brokenpenalty=10000\relax
Den weitaus größten Teil dieser Schrift nimmt eine Anleitung zur parallaktischen Berechnung und orthografischen Konstruktion von Sonnenfinsternissen ein. Sie erläutert das Verfahren praktisch an den Sonnenfinsternissen von 1816 und 1818 und gibt die Rechnungen für die erste in allen Einzelheiten an. Für völlige Anfänger kann diese Anleitung, die nur auf den gewöhnlichen einfachen Methoden beruht, von Nutzen sein. In dieser Hinsicht hätte der Verfasser bei den vorangestellten Vorschriften allerdings etwas mehr Sorgfalt auf eine klare und richtige Darstellung verwenden sollen. Auf den letzten 18 Seiten gibt der Verfasser noch die Hauptmomente der von 1818 bis 1860 in Brünn sichtbaren Finsternisse an, während die beigefügten Kupfertafeln eine allgemeine Darstellung derselben Finsternisse für die ganze Erde liefern. Darunter sind drei partielle, fünf ringförmige und vier totale Finsternisse. Für den Horizont von Brünn ist die Finsternis vom 8. Juli 1842 die größte, und die Größe der Verfinsterung beträgt dort 11 Zoll $13'$. Auf die Zeichnung selbst scheint nicht durchweg die größte Sorgfalt verwendet worden zu sein. So finden wir bei der erwähnten Finsternis die Punkte, an denen Anfang oder Ende der Finsternis mit Sonnenaufgang oder Sonnenuntergang zusammenfällt, als eine einzige zusammenhängende Linie eingezeichnet, obwohl sie doch notwendigerweise zwei getrennte Linien bilden mussten, die durch die Kontaktlinie zwischen dem südlichen Sonnenrand und dem nördlichen Mondrand verbunden sind. Die Rechnungen wurden nach \textsc{Triesnecker}s Mondtafeln in den Wiener Ephemeriden für 1803 und nach handschriftlichen Sonnentafeln desselben Astronomen ausgeführt. Der Nullmeridian ist anstelle der sonst üblichen $20^\circ$ auf $19^\circ\,53'\,45''$ westlich von Paris festgelegt.
\par\clearpage\endgroup
